%%%%%%%%%%%%%%%%%%%%%%%%%%%%%%%%%%%%%%%%%%%%%%%%%%%%%%%
% NOTAS TEORIA DE TRANSPORTE OPTIMO
%
% Contains ?? figures. 
% File started: ??
% First circulated version: ??
% Last version:
% Last modification: -
%%%%%%%%%%%%%%%%%%%%%%%%%%%%%%%%%%%%%%%%%%%%%%%%%%%%%%%

\documentclass[11pt,a4paper]{amsart}
\usepackage{xr-hyper}
\newcounter{chapter}

\usepackage{amssymb}
\usepackage{hyperref}
\usepackage{listings}
\usepackage{mathtools}
\usepackage{enumitem}

% Fuzz -------------------------------------------------------------------
\hfuzz2pt % Don't bother to report over-full boxes if over-edge is < 2pt
% Line spacing -----------------------------------------------------------

%%%%%%%%%fin pegada%%%%%%%%%%%%%%%%%%5

\usepackage[T1]{fontenc}
\usepackage[spanish]{babel}
\usepackage{graphicx}
\usepackage{tikz}
\usetikzlibrary{babel}

\setlength{\textwidth}{15.5cm} 
\setlength{\textheight}{24cm}
\setlength{\oddsidemargin}{0cm}
\setlength{\evensidemargin}{0cm}


\parskip 3pt

%%
%% OPERADORES
%%
\renewcommand{\Re}{\operatorname{Re}}
\newcommand{\tr}{\operatorname{tr}}
\newcommand{\id}{\operatorname{id}}
\def\div{\operatorname {\mathrm{div}}}
\def\Dom{\operatorname {\mathrm{Dom}}}
\def\argmin{\operatorname {\mathrm{arg\, min}}}
\def\argmax{\operatorname {\mathrm{arg\, max}}}
\def\gen{\operatorname {\mathrm{gen}}}
\def\Im{\operatorname {\mathrm{Im}}}
\def\Nu{\operatorname {\mathrm{Nu}}}
\def\sop{\operatorname {\mathrm{sop}}}
\def\diam{\operatorname {\mathrm{diam}}}
\def\dist{\operatorname {\mathrm{dist}}}
\def\var{\operatorname {\mathrm{Var}}}
\def\codim{\operatorname {\mathrm{codim}}}
\def\pint{\operatorname {--\!\!\!\!\!\int\!\!\!\!\!--}}

%%
%% EPSILON
%%
\newcommand{\ve}{\varepsilon}

%%
%% LETRAS MATHBB
%%
\newcommand{\R}{{\mathbb R}}
\newcommand{\Z}{{\mathbb Z}}
\newcommand{\C}{{\mathbb C}}
\newcommand{\Q}{{\mathbb Q}}
\newcommand{\N}{{\mathbb N}}
\newcommand{\E}{{\mathbb E}}
\newcommand{\V}{{\mathbb V}}
\renewcommand{\P}{{\mathbb P}}
\newcommand{\II}{{\mathbb I}}
\newcommand{\OO}{{\mathbb O}}
\newcommand{\Sb}{{\mathbb S}}
\newcommand{\Tb}{{\mathbb T}}

%%
%% LETRAS CALIFGRAFICAS
%%
\newcommand{\F}{{\mathcal F}}
\newcommand{\Hc}{{\mathcal H}}
\newcommand{\Tc}{{\mathcal T}}
\newcommand{\Ec}{{\mathcal E}}
\newcommand{\Vc}{{\mathcal V}}
\newcommand{\A}{{\mathcal A}}
\newcommand{\J}{{\mathcal J}}
\newcommand{\G}{{\mathcal G}}
\newcommand{\U}{{\mathcal U}}
\newcommand{\B}{{\mathcal B}}
\newcommand{\Sw}{{\mathcal S}}
\newcommand{\D}{{\mathcal D}}
\newcommand{\M}{{\mathcal M}}
\newcommand{\LL}{{\mathcal L}}

%%
%% LETRAS MATHBF
%%
\newcommand{\ab}{{\mathbf a}}
\newcommand{\bb}{{\mathbf b}}
\newcommand{\cc}{{\mathbf c}}
\newcommand{\ee}{{\mathbf e}}
\newcommand{\ff}{{\mathbf f}}
\newcommand{\gb}{{\mathbf g}}
\newcommand{\hh}{{\mathbf h}}
\newcommand{\n}{{\mathbf n}}
\newcommand{\mm}{{\mathbf m}}
\newcommand{\pp}{{\mathbf p}}
\newcommand{\qq}{{\mathbf q}}
\newcommand{\sbf}{{\mathbf s}}
\newcommand{\tb}{{\mathbf t}}
\newcommand{\uu}{{\mathbf u}}
\newcommand{\vv}{{\mathbf v}}
\newcommand{\ww}{{\mathbf w}}
\newcommand{\xx}{{\mathbf x}}
\newcommand{\yy}{{\mathbf y}}
\newcommand{\zz}{{\mathbf z}}
\newcommand{\I}{{\mathbf 1}}
\newcommand{\Pb}{{\mathbf P}}
\newcommand{\XX}{{\mathbf X}}
\newcommand{\YY}{{\mathbf Y}}

%%
%% LETRAS MATHFRAK
%%

\newcommand{\Gk}{{\mathfrak G}}


%%
%% TEOREMAS
%%
\newtheorem{teo}{Teorema}[section]
\newtheorem{lema}[teo]{Lema}
\newtheorem{prop}[teo]{Proposición}
\newtheorem{cor}[teo]{Corolario}

%%
%% REMARK
%%
\theoremstyle{remark}
\newtheorem{obs}[teo]{Observación}

%%
%% DEFINICION
%%
\theoremstyle{definition}
\newtheorem{defi}[teo]{Definición}
\newtheorem{ejem}[teo]{Ejemplo}
\newtheorem{ejer}[teo]{Ejercicio}
\newtheorem{nota}[teo]{Notación}
\newtheorem{ejercicio}{Ejercicio}[chapter]


%%
%% NUMERACION
%%
\numberwithin{subsection}{section}
\numberwithin{equation}{section}

%%
%% TITULO
%%


\externaldocument[N-]{../notas/Notas-TO}
\newcommand{\cuadernolink}[2]{\noindent\textbf{Cuaderno:} \emph{#1} (\href{https://colab.research.google.com/github/jfbonder/transporte-optimo/blob/main/cuadernos/#2}{abrir en Colab}).\par\medskip}
\pagestyle{plain}

\begin{document}
\renewcommand{\thechapter}{A}
\begin{center}
{\large\sc Teoría de Transporte Óptimo}\\[2pt]
{\small 2do cuatrimestre 2026 --- Julián Fernández Bonder}\\[10pt]
{\LARGE\bf Ejercicios del Apéndice A}\\[4pt]
{\large Complementos de teoría de la medida}\label{ap.medida}
\end{center}
\bigskip
\noindent{\small Los ejercicios son los de la sección de ejercicios del Apéndice~A de las notas del curso, con la misma numeración. Las referencias a teoremas, proposiciones y secciones remiten a las notas (\url{https://github.com/jfbonder/transporte-optimo}).}
\medskip


\begin{ejercicio}[Desintegraciones explícitas]\label{ejer.desint.ejemplos}
Hallar la desintegración $\{\pi_x\}$ respecto de la primera marginal de las
siguientes medidas en $\R^2$, y verificar la fórmula \eqref{N-eq:desint.int}
en cada caso:
\begin{enumerate}[label=(\alph*)]
\item $\pi=\mu\otimes\nu$ (producto);
\item $\pi=2\,\LL^2\lfloor_{\{0<y<x<1\}}$ (uniforme en un triángulo);
\item $\pi=(\id,T)_\#\LL^1\lfloor_{[0,1]}$ con $T(x)=x^2$;
\item $\pi=\frac12(\id,T_1)_\#\mu+\frac12(\id,T_2)_\#\mu$ con
$\mu=\LL^1\lfloor_{[0,1]}$, $T_1(x)=x$, $T_2(x)=1-x$ (un plan que no es un
mapa: ¿cuánto vale $\pi_x$?).
\end{enumerate}
\end{ejercicio}

\begin{ejercicio}[Independencia y planes inducidos]\label{ejer.desint.independencia}
Sea $\pi\in\Pi(\mu,\nu)$ con desintegración $\{\pi_x\}$.
\begin{enumerate}[label=(\alph*)]
\item Probar que $\pi=\mu\otimes\nu$ si y sólo si $\pi_x=\nu$ para
$\mu$-c.t.p.\ $x$.
\item Probar que $\pi$ está inducido por un mapa si y sólo si $\pi_x$ es una
masa de Dirac para $\mu$-c.t.p.\ $x$ (Proposición~\ref{N-prop.desint.mapa}), y
que en ese caso el mapa es $T(x)=\int y\,d\pi_x(y)$ (cuando $\nu$ tiene
primer momento finito).
\item Para $\pi\in\Pi(\mu,\nu)$ arbitrario con $\nu\in\mathcal P_2(\R^d)$,
sea $T(x)=\int y\,d\pi_x(y)$ el \emph{baricentro condicional}. Probar que
$\int|x-y|^2d\pi\ge\int|x-T(x)|^2d\mu$, con igualdad si y sólo si $\pi$ está
inducido por $T$. Concluir que si $T_\#\mu=\nu$, el plan $(\id,T)_\#\mu$ es
mejor que $\pi$: ``promediar el destino no aumenta el costo cuadrático''.
¿Por qué esto no demuestra que los planes óptimos son mapas?
\end{enumerate}
\end{ejercicio}

\begin{ejercicio}[Asociatividad del pegado]\label{ejer.pegado.asociativo}
Sean $\pi_{12}\in\Pi(\mu_1,\mu_2)$, $\pi_{23}\in\Pi(\mu_2,\mu_3)$ y
$\pi_{34}\in\Pi(\mu_3,\mu_4)$.
\begin{enumerate}[label=(\alph*)]
\item Probar que pegar primero $\pi_{12}$ con $\pi_{23}$ y el resultado con
$\pi_{34}$ (a lo largo de la tercera coordenada) da la misma medida en
$X_1\times X_2\times X_3\times X_4$ que pegar primero $\pi_{23}$ con
$\pi_{34}$ y luego con $\pi_{12}$.
\item Deducir la desigualdad triangular ``iterada''
\[
W_p(\mu_1,\mu_4)\le W_p(\mu_1,\mu_2)+W_p(\mu_2,\mu_3)+W_p(\mu_3,\mu_4)
\]
directamente a partir de un plan pegado, sin usar dos veces el
Teorema~\ref{N-teo.Wp.distancia}.
\item Probar que si los tres planes están inducidos por mapas $T_1,T_2,T_3$,
la medida pegada es $(\id,T_1,T_2\circ T_1,T_3\circ T_2\circ T_1)_\#\mu_1$.
\end{enumerate}
\end{ejercicio}

\begin{ejercicio}[Desintegración y transporte condicional]\label{ejer.desint.condicional}
Sean $\mu,\nu\in\mathcal P(\R^2)$ con marginales sobre la primera coordenada
$\mu^1$, $\nu^1$ y desintegraciones $\mu=\mu_x\otimes\mu^1$, $\nu=\nu_x\otimes\nu^1$.
\begin{enumerate}[label=(\alph*)]
\item (``Knothe--Rosenblatt''.) Supongamos $\mu^1$ y todas las $\mu_x$ sin
átomos. Sea $T^1$ el mapa monótono de $\mu^1$ a $\nu^1$ y, para cada $x$,
$T^2_x$ el mapa monótono de $\mu_x$ a $\nu_{T^1(x)}$. Probar que
$T(x,y)=(T^1(x),T^2_x(y))$ transporta $\mu$ en $\nu$. (Es el
\emph{reordenamiento de Knothe--Rosenblatt}; no es óptimo para el costo
cuadrático, pero es explícito y triangular.)
\item Probar que $W_2^2(\mu,\nu)\le W_2^2(\mu^1,\nu^1)+\int W_2^2(\mu_x,\nu_{T^1(x)})\,d\mu^1(x)$
y que, en general, la desigualdad es estricta.
\item Probar que si $\mu=\mu^1\otimes\alpha$ y $\nu=\nu^1\otimes\beta$ son
productos, entonces $W_2^2(\mu,\nu)=W_2^2(\mu^1,\nu^1)+W_2^2(\alpha,\beta)$.
\end{enumerate}
\end{ejercicio}

\end{document}
